\documentclass[12pt]{article}

% Настройка русского языка и шрифтов
\usepackage[utf8]{inputenc}
\usepackage[russian]{babel}
\usepackage[T2A]{fontenc}

% Математические пакеты
\usepackage{amsmath, amssymb, amsfonts}

% Настройка страницы
\usepackage{geometry}
\geometry{a4paper, margin=0.8in}

% Другие полезные пакеты
\usepackage{hyperref}
\usepackage{booktabs}
\usepackage{enumitem}
\usepackage{longtable}
\usepackage{graphicx}
\usepackage{xcolor}

% Определение стилей для комментариев
\definecolor{commentgray}{rgb}{0.4,0.4,0.4}
\definecolor{antigravityblue}{rgb}{0.0, 0.45, 0.74}

\newcommand{\commentary}[1]{{\small\itshape\color{commentgray} #1}}
\newcommand{\antigravity}[1]{{\small\color{antigravityblue}\textbf{Antigravity:} #1}}

% Настройка ссылок
\hypersetup{
    colorlinks=true,
    linkcolor=blue,
    filecolor=magenta,      
    urlcolor=cyan,
    pdftitle={Глубокий разбор: Линейные модели},
}

\title{Глубокий разбор: Линейные модели в машинном обучении}
\author{Твой персональный гид по ML}
\date{\today}

\begin{document}

\subsection{Типология задач машинного обучения}
В обучении с учителем выделяют следующие задачи:
\begin{itemize}
    \item \textbf{Регрессия:} целевая переменная принимает вещественные значения ($Y = \mathbb{R}$).
    \item \textbf{Бинарная классификация:} выбор из двух классов, обычно $\{0, 1\}$ или $\{-1, 1\}$.
    \item \textbf{Многоклассовая классификация:} выбор одного из $K$ классов.
    \item \textbf{Многоклассовая классификация multilabel:} объект может принадлежать нескольким классам одновременно ($Y = \{0, 1\}^K$).
    \item \textbf{Ранжирование:} необходимо отсортировать объекты по их релевантности.
\end{itemize}

\commentary{Главное отличие в природе таргетов. В регрессии мы предсказываем число, в классификации — метку класса.}

\subsection{Математический аппарат и обозначения}
\begin{itemize}
    \item \textbf{Матрица $X$}: матрица «объекты — признаки» размера $N \times D$.
    \item \textbf{Вектор $y$}: истинные ответы.
    \item \textbf{Вектор весов $w$}: параметры модели. $y = \langle x, w \rangle + w_0$.
    \item \textbf{Решение МНК}: $w = (X^TX)^{-1}X^Ty$.
    \item \textbf{Численные методы}: Для стабильного обращения матриц используют QR-разложение ($X=QR \implies w = R^{-1}Q^Ty$) или сингулярное разложение (SVD: $X=U\Sigma V^T \implies w = V \Sigma^{-1} U^T y$).
\end{itemize}

\subsection{Обработка категориальных признаков}
\textbf{One-hot-кодирование} заменяет категориальный признак с $M$ значениями на $M$ новых бинарных признаков.

\textbf{Почему нельзя (1, 2, 3):} Линейная модель воспринимает эти числа как шкалу. Она решит, что категория «3» в три раза больше «1», что бессмысленно для профессий или городов.

\subsection{Интерпретируемость моделей}
Это способность модели объяснить, на основании чего сделано предсказание.
\begin{itemize}
    \item \textbf{Линейные модели:} через веса $w_i$. Большой положительный вес — признак сильно увеличивает результат.
    \item \textbf{Решающие деревья:} через структуру узлов — путь от корня к листу.
\end{itemize}

\subsection{Свойства и ограничения линейных моделей}
\begin{itemize}
    \item \textbf{Узкий класс:} описывают только линейные зависимости. Требуют \textit{feature engineering}.
    \item \textbf{Мультиколлинеарность:} зависимость признаков ведет к огромным и неустойчивым весам.
    \item \textbf{Масштаб:} признаки должны быть отнормированы для правильной интерпретации весов.
\end{itemize}

\subsection{Сведение обучения к задаче оптимизации}
\begin{itemize}
    \item \textbf{Функция потерь (Loss)}: ошибка на одном объекте.
    \item \textbf{Функционал качества}: средняя ошибка по выборке (Эмпирический риск).
    \item \textbf{SGD (Stochastic Gradient Descent)}: Оценка градиента на подвыборке — батче (batch). Обновление весов: $w_{next} = w - \eta \cdot \frac{1}{B} \sum_{t=1}^{B} \nabla_{w}L(w,x_{i_t},y_{i_t})$. Это на порядки быстрее классического градиентного спуска на больших данных.
\end{itemize}

\subsection{Мультиколлинеарность и регуляризация}
\textbf{Регуляризация} добавляет штраф за сложность: $\min (Loss + \lambda \|w\|)$.
Коэффициент $\lambda$ контролирует баланс между ошибкой и величиной весов.

\subsection{Сравнение L1 и L2 регуляризаций}
\begin{itemize}
    \item \textbf{$L_2$ (Ridge):} $\lambda \sum w_i^2$. Делает веса маленькими, устойчива к шуму.
    \item \textbf{$L_1$ (Lasso):} $\lambda \sum |w_i|$. Обнуляет веса неважных признаков (разреживание).
\end{itemize}

\subsection{Функция потерь перцептрона}
$F(M) = \max(0, -M)$, где $M = y_i \langle w, x_i \rangle$ — отступ. Ошибка растет, если отступ отрицательный.

\subsection{Hinge Loss и SVM}
$H(M) = \max(0, 1 - M)$. Заставляет модель максимизировать \textbf{зазор} (margin). Опорные векторы лежат на границе этого зазора.

\subsection{Логистическая регрессия и вероятности}
Перевод логита $\langle w, x \rangle$ в вероятность через сигмоиду: $\sigma(z) = \frac{1}{1 + e^{-z}}$.

\subsection{Анализ гиперпараметров регуляризации}
\begin{itemize}
    \item \textbf{Большой $\lambda$:} веса стремятся к нулю (недообучение).
    \item \textbf{Маленький $\lambda$:} сложная граница, подстройка под выбросы (переобучение).
\end{itemize}

\subsection{Softmax для многоклассовой классификации}
$P(y=k|x) = \frac{e^{z_k}}{\sum e^{z_j}}$. Гарантирует, что сумма вероятностей всех классов равна 1.

\vspace{1cm}
\antigravity{Линейные модели — это фундамент. Правильно настроенная регуляризация и грамотный фича-инжиниринг делают их мощным инструментом даже в эпоху глубокого обучения.}

\end{document}
